\documentclass[10pt, xcolor={dvipsnames}]{beamer}

%========================
% CẤU HÌNH HỆ THỐNG
%========================
\usepackage[utf8]{inputenc}
\usepackage[T5]{fontenc}
\usepackage[vietnamese]{babel}
\usepackage{tikz}
\usetikzlibrary{arrows.meta, positioning, calc, shapes.geometric, shadows, patterns}
\usepackage{amsmath, amssymb, amsthm}
\usepackage{tcolorbox}
\tcbuselibrary{skins, shadows, themes}

%========================
% TÔNG MÀU & STYLE
%========================
\usetheme{Madrid}
\definecolor{mainblue}{RGB}{0, 102, 104}
\definecolor{darkblue}{RGB}{0, 32, 96}
\definecolor{lightorange}{RGB}{255, 140, 0}
\definecolor{softblue}{RGB}{235, 245, 255}

\setbeamercolor{palette primary}{bg=mainblue, fg=white}
\setbeamercolor{palette secondary}{bg=lightorange, fg=white}
\setbeamercolor{structure}{fg=mainblue}

% Hộp nội dung đặc sắc
\newtcolorbox{detailbox}[2]{
    enhanced, colback=white, colframe=#1, title=#2,
    fonttitle=\bfseries, arc=3pt, boxrule=1pt, shadow={1mm}{-1mm}{0mm}{black!10}
}

\title{LÝ THUYẾT ĐỒ THỊ}
\subtitle{Toàn tập: Từ Cơ bản đến Nâng cao}
\author{Nguyễn Như Bảo}
\date{Học kỳ II - 2026}

\begin{document}

\frame{\titlepage}

%===========================================================
% PHẦN 1: GIỚI THIỆU (10 SLIDE)
%===========================================================
\section{Phần 1: Giới thiệu tổng quan}

\begin{frame}{1.1. Bài toán khởi nguồn: Bảy cây cầu Königsberg}
    \begin{columns}
        \begin{column}{0.5\textwidth}
            \begin{detailbox}{mainblue}{Lịch sử (1736)}
                Leonhard Euler đã chứng minh không thể đi qua cả 7 cây cầu mà mỗi cầu chỉ đi đúng 1 lần.
            \end{detailbox}
            \begin{itemize}
                \item \textbf{Đỉnh:} 4 vùng đất.
                \item \textbf{Cạnh:} 7 cây cầu nối.
                \item \textbf{Phát hiện:} Hình dáng địa lý không quan trọng bằng kết nối.
            \end{itemize}
        \end{column}
        \begin{column}{0.5\textwidth}
            \centering
            \begin{tikzpicture}[node distance=1.2cm, every node/.style={circle, draw=mainblue, fill=blue!10, thick}]
                \node (A) {A}; \node (B) [above right=of A] {B};
                \node (C) [below right=of A] {C}; \node (D) [right=of A, xshift=1.5cm] {D};
                \draw (A) to [bend left=20] (B); \draw (A) to [bend right=20] (B);
                \draw (A) to [bend left=20] (C); \draw (A) to [bend right=20] (C);
                \draw (B) -- (D); \draw (C) -- (D); \draw (B) -- (C);
            \end{tikzpicture}
        \end{column}
    \end{columns}
\end{frame}

\begin{frame}{1.2. Định nghĩa hình thức $G = (V, E)$}
    \begin{detailbox}{darkblue}{Cấu trúc cốt lõi}
        Đồ thị $G$ được xác định bởi tập đỉnh $V$ và tập cạnh $E$.
    \end{detailbox}
    \begin{itemize}
        \item \textbf{V (Vertices):} Tập các đối tượng hữu hạn.
        \item \textbf{E (Edges):} Tập các cặp đỉnh $(u, v)$ thể hiện quan hệ.
    \end{itemize}
    \centering
    \begin{tikzpicture}[scale=0.8]
        \node[draw,circle](1) at (0,1){1}; \node[draw,circle](2) at (2,1){2};
        \node[draw,circle](3) at (1,0){3}; \draw (1)--(2)--(3)--(1);
    \end{tikzpicture}
    \\ \small $V = \{1, 2, 3\}$, $E = \{(1,2), (2,3), (3,1)\}$
\end{frame}

\begin{frame}{1.3. Ứng dụng: Mạng xã hội (Social Graphs)}
    \begin{columns}
        \begin{column}{0.6\textwidth}
            \begin{itemize}
                \item \textbf{Đỉnh:} Người dùng Facebook.
                \item \textbf{Cạnh:} Quan hệ "Bạn bè".
                \item \textbf{Ứng dụng:} Thuật toán gợi ý người quen (Triangle Closure).
            \end{itemize}
            \begin{exampleblock}{Tư duy}
                Nếu A bạn B, B bạn C, khả năng cao A và C sẽ biết nhau.
            \end{exampleblock}
        \end{column}
        \begin{column}{0.4\textwidth}
            \centering
            \begin{tikzpicture}
                \node[draw,circle,fill=blue!20](A){A}; \node[draw,circle,fill=blue!20,right=of A](B){B};
                \node[draw,circle,fill=blue!20,below=of B](C){C};
                \draw (A)--(B); \draw (B)--(C); \draw[dashed, red] (A)--(C);
            \end{tikzpicture}
        \end{column}
    \end{columns}
\end{frame}

\begin{frame}{1.4. Ứng dụng: Google Maps & Điều hướng}
    \begin{detailbox}{lightorange}{Hệ thống Giao thông}
        Mô hình hóa bản đồ thế giới thành các nút và đường đi.
    \end{detailbox}
    \begin{itemize}
        \item \textbf{Đỉnh:} Giao lộ, thành phố.
        \item \textbf{Trọng số cạnh:} Chiều dài đường, thời gian kẹt xe.
        \item \textbf{Bài toán:} Tìm đường đi ngắn nhất (Dijkstra Algorithm).
    \end{itemize}
\end{frame}

\begin{frame}{1.5. Ứng dụng: World Wide Web (Web Graph)}
    \begin{columns}
        \begin{column}{0.5\textwidth}
            \centering
            \begin{tikzpicture}[node distance=1cm]
                \node[draw,rectangle](W1){URL A}; \node[draw,rectangle,below=of W1](W2){URL B};
                \draw[->,>=Stealth] (W1) -- (W2);
            \end{tikzpicture}
        \end{column}
        \begin{column}{0.5\textwidth}
            \begin{itemize}
                \item \textbf{Đỉnh:} Trang Web.
                \item \textbf{Cạnh:} Siêu liên kết (Link).
                \item \textbf{Bản chất:} Đồ thị có hướng khổng lồ.
            \end{itemize}
        \end{column}
    \end{columns}
\end{frame}

\begin{frame}{1.6. Ứng dụng: Hóa học & Sinh học}
    \begin{itemize}
        \item \textbf{Hóa học:} Cấu trúc phân tử (Đỉnh = Nguyên tử, Cạnh = Liên kết hóa học).
        \item \textbf{Sinh học:} Lưới thức ăn, tương tác Protein.
    \end{itemize}
    \centering
    \begin{tikzpicture}[scale=0.6]
        \node[draw,circle](C) at (0,0){C}; \foreach \a in {0,90,180,270} \node[draw,circle] at (\a:1.5){H};
        \foreach \a in {0,90,180,270} \draw (C) -- (\a:1.2);
    \end{tikzpicture}
\end{frame}

\begin{frame}{1.7. Ứng dụng: Mạng máy tính & Viễn thông}
    \begin{detailbox}{mainblue}{Cơ sở hạ tầng}
        Đảm bảo gói tin đi từ nguồn đến đích hiệu quả nhất.
    \end{detailbox}
    \begin{itemize}
        \item \textbf{Đỉnh:} Router, Switch, Máy chủ.
        \item \textbf{Cạnh:} Cáp quang, kết nối không dây.
    \end{itemize}
\end{frame}

\begin{frame}{1.8. Ứng dụng: Phân tích tài chính & Giao dịch}
    \begin{itemize}
        \item \textbf{Đỉnh:} Tài khoản ngân hàng.
        \item \textbf{Cạnh:} Lệnh chuyển tiền.
        \item \textbf{Mục đích:} Phát hiện rửa tiền thông qua các chu trình (cycles) bất thường.
    \end{itemize}
\end{frame}

\begin{frame}{1.9. Đồ thị trong Trí tuệ nhân tạo (AI)}
    \begin{detailbox}{darkblue}{Graph Neural Networks (GNN)}
        Xu hướng mới trong Deep Learning: Học trực tiếp trên cấu trúc đồ thị để dự đoán tính chất dữ liệu.
    \end{detailbox}
\end{frame}

\begin{frame}{1.10. Tổng kết Phần 1}
    \begin{enumerate}
        \item Đồ thị là công cụ mô hình hóa vạn năng.
        \item Xuất hiện từ MXH đến khoa học cốt lõi.
        \item Hiểu đồ thị là bước đầu để làm chủ giải thuật phức tạp.
    \end{enumerate}
\end{frame}


\end{document}
